% Generated draft: requires mathematical and source review.
\documentclass[11pt]{article}
\usepackage[margin=25mm]{geometry}
\usepackage{fontspec}
\setmainfont{Latin Modern Roman}
\usepackage{amsmath,amssymb,mathtools,unicode-math}
\setmathfont{Latin Modern Math}
\usepackage{enumitem,needspace,xurl,hyperref}
\setlength{\parindent}{0pt}
\setlength{\parskip}{5pt}
\setlength{\emergencystretch}{3em}
\newcommand{\N}{\mathbb{N}}
\newcommand{\Z}{\mathbb{Z}}
\newcommand{\Q}{\mathbb{Q}}
\newcommand{\R}{\mathbb{R}}
\newcommand{\C}{\mathbb{C}}
\newcommand{\F}{\mathbb{F}}
\newcommand{\A}{\mathbb{A}}
\newcommand{\PP}{\mathbb{P}}
\newcommand{\E}{\mathbb{E}}
\newcommand{\eps}{\varepsilon}
\newcommand{\dd}{\,\mathrm{d}}
\begin{document}
\section*{UnsolvedMath ID: 30001675}
\textbf{Source problem code:} catalogue entry
\subsection{Definitions and mathematical statement}
Let \(\N=\{0,1,2,\ldots\}\), and let \([\N]^\omega\) denote its infinite subsets, identified with their characteristic functions and equipped with the subspace topology inherited from \(2^{\N}\). An infinite family \(\mathcal A\subseteq[\N]^\omega\) is a maximal almost disjoint family (a MAD family) if distinct members have finite intersection and
\[
\forall X\in[\N]^\omega\ \exists A\in\mathcal A\quad |X\cap A|=\aleph_0.
\]
Thus maximality is with respect to adding infinite subsets of \(\N\).

For \(f,g\in\N^{\N}\), write \(f\leq^*g\) if \(f(n)\leq g(n)\) for all sufficiently large \(n\). The bounding number is
\[
\mathfrak b=\min\bigl\{|B|:B\subseteq\N^{\N}\text{ and }\forall g\in\N^{\N}\ \exists f\in B\ (f\not\leq^*g)\bigr\}.
\]

The supplied record does not define \(a_{\mathrm{Borel}}\) or \(a_{\mathrm{closed}}\). Their intended minimization problems and ambient topological spaces must therefore be recovered from the cited source before these parts can receive a complete mathematical formulation. In particular, this draft does not identify them with the cardinalities of Borel or closed MAD families.

The projective notation also requires a convention about real parameters. Under the boldface interpretation, a \(\Pi^1_1\) subset of \([\N]^\omega\) is coanalytic, meaning that its complement is the projection of a Borel subset of \([\N]^\omega\times\N^{\N}\). A \(\Sigma^1_2\) subset is the projection of a coanalytic subset of that product. These definitions allow real parameters. The record does not explicitly distinguish this interpretation from the lightface, effective interpretation. For either fixed interpretation \(\Gamma\), let
\[
P_{\Gamma}\quad\Longleftrightarrow\quad
\exists\mathcal A\subseteq[\N]^\omega\ 
\bigl(\mathcal A\text{ is a MAD family and }\mathcal A\in\Gamma\bigr).
\]
The same parameter convention must be used throughout.

For a set \(S\subseteq[\N]^\omega\), use the homogeneous-set formulation of the Ramsey property:
\[
\forall X\in[\N]^\omega\ \exists Y\in[X]^\omega\quad
\bigl([Y]^\omega\subseteq S\ \lor\ [Y]^\omega\cap S=\varnothing\bigr).
\]
Write \(R_{\Sigma^1_2}\) for the assertion that every \(\Sigma^1_2\) subset of \([\N]^\omega\) has this property. This states the convention used here; the source record does not specify its precise Ramsey-property convention.

The supplied problem bundles the following questions.
\begin{enumerate}
\item Determine the intended equality or consistency question involving
\[
a_{\mathrm{Borel}}=a_{\mathrm{closed}},
\]
and the consistency question involving
\[
a_{\mathrm{Borel}}>\mathfrak b.
\]
The original wording, ``Are \(a_{\mathrm{Borel}}=a_{\mathrm{closed}}\) and is \(a_{\mathrm{Borel}}>\mathfrak b\) consistent?'', can ask whether the equality holds and whether the strict inequality is consistent. The cleaned wording can instead ask consistency of each assertion, or of their conjunction. These readings are retained as unresolved alternatives.
\item Is it consistent that
\[
\mathfrak b>\aleph_1\quad\text{and}\quad P_{\Pi^1_1}?
\]
\item Is the existence equivalence
\[
P_{\Pi^1_1}\ \Longleftrightarrow\ P_{\Sigma^1_2}
\]
valid?
\item Does
\[
R_{\Sigma^1_2}\ \Longrightarrow\ \neg P_{\Sigma^1_2}
\]
hold?
\end{enumerate}
The background theory for consistency and validity is not stated in the record. Under the usual ZFC reading, question 2 asks about consistency of \(\mathrm{ZFC}+(\mathfrak b>\aleph_1)+P_{\Pi^1_1}\), and questions 3 and 4 ask whether their displayed implications or equivalence hold in every model of ZFC, with the chosen projective conventions. This is an explicit conditional interpretation, not an additional hypothesis supplied by the source. The undefined cardinal invariants and the alternative readings in question 1 remain unresolved.
\subsection{Short English statement}
The problem asks about comparisons between two unspecified definable-MAD invariants and the bounding number, the coexistence of a large bounding number with a coanalytic MAD family, equivalence of two projective MAD existence assertions, and whether projective Ramsey regularity rules out projective MAD families. The invariant definitions, part of the first question, and the logical and definability conventions require source clarification. The supplied status is partially solved; this draft makes no independent claim about which questions remain open.
\subsection{Sources}
\begin{itemize}
\item Local registry: \texttt{data/unsolvedmath/problems.json}, permanent UnsolvedMath ID \(30001675\), title \emph{Definable Maximal Almost Disjoint Families}, exact source code \texttt{catalogue entry}; supplied collection: Oberwolfach Reports Open Problems. The recorded proposed year is 2011; no individual proposer is supplied.
\item Supplied primary provenance: \emph{Set Theory} (2011), Oberwolfach Reports, pp.\ 85--140, DOI \texttt{10.4171/owr/2011/02}. The report documents workshop discussions and open problems. Supplied source URL and literature-source label ``Oberwolfach problem statement'': \url{https://doi.org/10.4171/owr/2011/02}.
\item Supplied literature source, labeled ``Tornquist, projective MAD families'': \url{https://doi.org/10.2178/jsl.7804080}. No fuller bibliographic information is supplied.
\item Supplied literature source, labeled ``strong projective witnesses'': \url{https://doi.org/10.1016/j.apal.2022.103135}. No author or fuller bibliographic information is supplied.
\item Literature assessment dated exactly \texttt{2026-08-21}: the supplied curation labels the bundle ``partially solved'' and reports mixed status, with several projective MAD implications and consistency results known, while the equality and strict-inequality questions remain unresolved. This is supplied dated curation, not independently verified here; it does not identify precisely which individual projective assertions have been settled. The source expressly requests independent verification of time-sensitive status and cited claims.
\item The supplied statement-status label is \texttt{corrected\_verified}. Its accompanying verification note says only that TeX notation was enclosed in math delimiters without changing mathematical content; this does not establish the missing definitions or verify the mathematical assertions.
\end{itemize}
\end{document}
